\section{Dataset Construction}
\label{sec:dataset-construction}

Figure~\ref{fig:pipeline} illustrates how each multi-turn case is progressively
downgraded from a successful strong-model trajectory to produce per-step
verified tier labels through execution-verified search.

\begin{figure}[t]
\centering
\includegraphics[width=\linewidth]{figures/fig2.pdf}
\caption{TwinRouterBench construction pipeline. For each multi-turn case, the
pipeline starts from a successful strong-model trajectory and progressively
downgrades individual steps to cheaper tiers via execution-verified search,
producing the verified tier label for every LLM call in the trace.}
\label{fig:pipeline}
\end{figure}

\paragraph{Seed from successful strong trajectories.}
We first collect trajectories that resolve under a strong-model setting. For
SWE-bench, Claude Opus 4.6 or GPT-5.4 running through
\texttt{mini-swe-agent} must submit a patch that passes the official
\texttt{FAIL\_TO\_PASS} tests; BFCL, mtRAG, QMSum, and PinchBench~\citep{pinchbench} use their
respective harness pass predicates. Failed seeds are dropped so that routing
labels do not conflate task difficulty with routing error.


\paragraph{Search lower tiers under causal prefixes.}
For each surviving trajectory, Claude Opus 4.6 provides a search hint:
whether a step appears downgradeable and, if so, the most aggressive lower tier
$h_i$ to probe first. The hint is a pruning device, never a label. We then run
a greedy sequential-locking downgrade search (Appendix Algorithm~\ref{alg:seq-lock}).
At step $i$, previously locked steps keep their verified tiers, future steps
use the strong tier, and candidates for step $i$ are tried from cheapest to
strongest starting at $h_i$. Because per-step correctness cannot always be
observed directly, we accept a candidate when the full trajectory still passes
with the same number of steps, inferring that the downgraded step was handled
correctly; if no candidate passes, the step stays at the strong tier. This
reduces the naive $|T|^N$ grid to $\mathcal{O}(|T|N)$ trials while ensuring
later prefixes contain actual degraded outputs from earlier routed steps.
Locked prefixes are replayed from a response cache keyed by instance, step,
model, and generation parameters; cache entries are invalidated when an
upstream lock changes.

\paragraph{Verify tiers as pools, not single models.}
A tier label should mean that the tier pool can solve the step, not that one
model happened to pass. We run a three-model cascade for non-\texttt{low}
labels: a tier is accepted if any of the three pool models passes under the
mixed-model trace. In a SWE-bench last-step audit, the cascade recovered 28 of
33 downgradeable labels versus 15 of 33 under single-model probing, confirming
that single-representative search is too brittle for supervision.

\paragraph{Harden open-ended pass predicates.}
Some benchmarks (mtRAG, QMSum) rely on LLM-as-judge evaluation whose default
prompts and rubrics are not strict enough. During early development we manually
inspected a random sample of judge verdicts and found pseudo-passes where
surface-level similarity masks incorrect or incomplete answers. We replace these
weak reference-similarity rubrics with a structured judge that first checks
whether the current turn is solved, then scores Faithfulness, Appropriateness,
and Completeness, with evidence-conflict hard failures and a conservative
fail-on-uncertainty default. We calibrate the tightened rubric against
Radbench-style thresholded ratings \citep{kuo2025radbench} and verify through
iterative test runs that the hardened judge eliminates the pseudo-pass patterns
discovered in our initial sample; structurally
unreliable mtRAG samples and all-tier failures are discarded. SWE-bench and
BFCL already have executable or structural pass predicates and need no
additional hardening.

\paragraph{Manual audit.}
Our downgrade search fixes all other steps and tests one step at a time. This
greedy strategy reduces computational cost from exponential to linear but may
miss cross-step interaction effects: a combination of two downgrades might fail
even though each succeeds individually. To guard against such cases, we conduct
a human sampling audit over the three workloads with multi-step trajectories
(BFCL, SWE-bench, PinchBench), drawing an independent ${\sim}10\%$ random
sample of routed steps from each. mtRAG and QMSum are single-turn tasks that
lack multi-turn compounding complexity; their labels are already constrained by
the hardened judge whose pseudo-pass patterns were eliminated through the
iterative calibration described above, so they are not resampled in the final
audit. For each sampled step a reviewer freezes the router-visible prefix
and judges whether the tier can be lowered one more level, labelling it as
\emph{tight}, \emph{uncertain}, or \emph{further-downgradeable}. The audit
covers 64 steps (25 BFCL, 34 SWE-bench, 5 PinchBench): 63 were judged tight;
one SWE-bench step was flagged as further-downgradeable and its tier was lowered
before release (Table~\ref{tab:manual-audit-main}). Details of the audit
protocol are in Appendix~\ref{app:manual-audit}.

\begin{table}[t]
\centering
\small
\begin{adjustbox}{max width=\linewidth}
\begin{tabular}{lrrrrr}
\toprule
Benchmark & Total steps & Sampled steps & Sampling \% & Downgradeable & Tight \% \\
\midrule
BFCL & 248 & 25 & 10.1 & 0 & 100.0 \\
SWE-bench & 336 & 34 & 10.1 & 1 & 97.1 \\
PinchBench & 48 & 5 & 10.4 & 0 & 100.0 \\
\midrule
Total & 632 & 64 & 10.1 & 1 & 98.4 \\
\bottomrule
\end{tabular}
\end{adjustbox}
\caption{Step-level manual audit summary. \emph{Downgradeable}: the reviewer
believes the verified tier can be lowered one more level without breaking task
resolution. \emph{Tight~\%}: the share already at the conditional optimum.
We sample routed steps from the three workloads with multi-step trajectories;
mtRAG and QMSum are single-turn tasks already guarded by the hardened judge and
are not resampled.}
\label{tab:manual-audit-main}
\end{table}

\paragraph{Release criterion.}
No row enters the static track unless (1)~the strong seed trajectory passed,
(2)~the tier pool contains a model that passes under the mixed trace, and
(3)~for open-ended workloads (mtRAG, QMSum), the hardened judge does not flag
the answer as a pseudo-pass. Each label is
an execution-verified estimate of the cheapest tier that can correctly handle
the current step: since direct per-step judgement is not always possible, we
infer per-step sufficiency from end-to-end trajectory pass with a fixed step
count, and cross-check through manual audit.
